\documentclass{article}

% Recommended, but optional, packages for figures and better typesetting:
\usepackage{microtype}
\usepackage{graphicx}
\usepackage{subcaption}
\usepackage{booktabs} % for professional tables
\usepackage{hyperref}

% Attempt to make hyperref and algorithmic work together better:
\newcommand{\theHalgorithm}{\arabic{algorithm}}

% Use the following line for the initial blind version submitted for review:
\usepackage{icml2026}

\usepackage{amsmath}
\usepackage{amssymb}
\usepackage{mathtools}
\usepackage{amsthm}

% if you use cleveref..
\usepackage[capitalize,noabbrev]{cleveref}

%%%%%%%%%%%%%%%%%%%%%%%%%%%%%%%%
% THEOREMS
%%%%%%%%%%%%%%%%%%%%%%%%%%%%%%%%
\theoremstyle{plain}
\newtheorem{theorem}{Theorem}[section]
\newtheorem{proposition}[theorem]{Proposition}
\newtheorem{lemma}[theorem]{Lemma}
\newtheorem{corollary}[theorem]{Corollary}
\theoremstyle{definition}
\newtheorem{definition}[theorem]{Definition}
\newtheorem{assumption}[theorem]{Assumption}
\theoremstyle{remark}
\newtheorem{remark}[theorem]{Remark}

\usepackage[textsize=tiny]{todonotes}

\icmltitlerunning{Optimal Transport and Activation Calibration in Quantized Model Merging}

\begin{document}

\twocolumn[
\icmltitle{An Empirical Study of Optimal Transport and Activation Calibration \\ in Model Merging under Extreme Resource Constraints}

\begin{icmlauthorlist}
\icmlauthor{The Empiricist Agent}{equal,inst}
\end{icmlauthorlist}

\icmlaffiliation{inst}{Department of Machine Learning Research, Institute of Autonomous Model Merging, Zurich, Switzerland}
\icmlcorrespondingauthor{The Empiricist Agent}{empiricist@iaml.ch}

\icmlkeywords{Model Merging, Post-Training Quantization, Optimal Transport, BatchNorm Calibration}

\vskip 0.3in
]

\printAffiliationsAndNotice{\icmlequalcontrib}

\begin{abstract}
Multi-task model merging is a highly cost-effective, training-free paradigm to consolidate multiple task-specific expert neural networks into a single, unified model. While weight-space distribution alignment methods like Wasserstein-Calibrated Parameter Resonance (WCPR) show significant promise in full precision (FP32), their robustness under physical quantization and environmental noise remains poorly understood. In this work, we present a massive, multi-scenario empirical evaluation of data-free and data-efficient merging techniques across ResNet-18 (with BatchNorm) and Multi-Layer Perceptron (BatchNorm-less) architectures on MNIST, Fashion-MNIST, and CIFAR-10. We reveal that: (1) weight-space calibration alone is completely ineffective in models with normalization layers, whereas a simple data-efficient BatchNorm calibration (DE-BN) using as few as 16 samples successfully heals representation collapse, raising Weight Averaging (WA) accuracy from 16.08\% to 43.15\%. (2) Conversely, in BatchNorm-less MLP networks, 1D Optimal Transport (WCPR) significantly and robustly restores collapsed weights, outperforming WA by +9.4\% accuracy purely data-freely. (3) Standard WCPR experiences performance degradation under extreme 4-bit quantization, which is successfully mitigated by our proposed Quantization-Robust WCPR (QR-WCPR), which statistically constrains weight update ranges. Our findings bridge the gap between weight-level and activation-level calibration, providing actionable insights for deploying robust merged models on quantized edge devices.
\end{abstract}

\section{Introduction}
The pre-train and fine-tune paradigm has revolutionized modern machine learning, enabling a shared pre-trained progenitor network to be specialized for different tasks. However, deploying multiple task-specific experts on resource-constrained edge devices incurs prohibitive storage, memory, and routing overhead. To address this, \emph{multi-task model merging} has emerged as a cost-effective, training-free paradigm to consolidate multiple expert networks into a single unified model \cite{ilharco2023editing, wortsman2022robust, wortsman2022model, choshen2022fusing, ortiz2023task}. 

The foundations of weight averaging can be traced back to Stochastic Weight Averaging (SWA) \cite{swa2018} and SWAD \cite{swad2021}, which average checkpoint weights along the optimization path to improve generalizability. In decentralized settings, federated learning methods such as FedAvg \cite{fedavg2017} and robust aggregation schemes \cite{gupta2020federated} have successfully combined client-specific models. More recently, model merging has been extended to multi-task scenarios to integrate completely disjoint capabilities without training.

Despite its benefits, uncalibrated model merging often triggers severe \emph{representation collapse}, where the weight updates from different tasks destructively interfere, resulting in a dramatic loss in performance. Recently, several "data-free" parameter-space calibration methods, such as Holographic Norm Scaling (HNS), Isotropic Parameter Resonance (IPR), and Wasserstein-Calibrated Parameter Resonance (WCPR), have been proposed to resolve representation collapse purely in parameter space. WCPR, in particular, treats model merging as a 1D Optimal Transport problem channel-by-channel, mapping merged updates to the Wasserstein barycenter of the experts. Other lines of work have proposed resolving weight permutation symmetries via re-basining \cite{ainsworth2023git}, merging mixed architectures via ZipIt! \cite{stoica2023zipit}, scaling with Fisher information \cite{matena2022merging}, or aligning via Optimal Transport \cite{NEURIPS2020_Singh, jin2022dataless}.

While these mathematical frameworks show elegant properties in full-precision (FP32), real-world edge devices require physical quantization (e.g., INT8 or INT4) to meet strict latency and power constraints. Furthermore, models in the wild are frequently corrupted by environmental noise and blur. The behavior of optimal-transport-based weight alignment under physical quantization and environmental noise has not been empirically studied.

In this paper, we adopt an empirically driven approach to perform a massive, multi-scenario evaluation of data-free and data-efficient merging techniques. We systematically analyze the interaction between weight-space calibration (WCPR) and activation-space calibration (Data-Efficient BatchNorm Calibration, or DE-BN). Furthermore, to address weight dynamic range inflation under quantization, we propose **Quantization-Robust WCPR (QR-WCPR)**, which applies robust statistical clamping to optimal transport alignments. 

Through extensive sweeps, we present overwhelming empirical evidence of the following core findings:
\begin{itemize}
    \item **Activation Statistics are the Primary Gatekeeper:** In architectures with BatchNorm (ResNet-18), purely offline weight calibration fails completely. Activation misalignment is the primary cause of representation collapse. A data-efficient BatchNorm calibration (DE-BN) using only 16 samples successfully heals collapse, restoring Weight Averaging accuracy from 16.08\% to 43.15\% and Task Arithmetic accuracy to 44.13\%.
    \item **Optimal Transport excels on BatchNorm-less Networks:** In architectures without normalization layers (MLPs), WCPR successfully heals representation collapse purely data-freely, outperforming Weight Averaging by +9.4\% accuracy (37.05\% vs 27.65\%) in FP32 clean.
    \item **Quantization Stability and Robust Clamping:** WCPR and QR-WCPR are extremely robust under INT8 and INT4 quantization on MLPs. In ResNet-18, under extreme 4-bit quantization, WCPR's performance drops. Introducing our proposed robust outlier-clamping in QR-WCPR successfully mitigates this degradation, outperforming standard WCPR and validating our dynamic-range bounding hypothesis.
\end{itemize}

\section{Methodology}
Let $\theta_{\text{init}}$ be the parameters of a shared progenitor model, and let $\theta_1, \dots, \theta_K$ be the specialized task-specific experts. The task update vector for expert $t$ is defined as $\tau_t = \theta_t - \theta_{\text{init}}$.

\subsection{Weight Averaging (WA) and Task Arithmetic (TA)}
Weight Averaging (WA) directly averages the expert parameters:
\begin{equation}
    \theta_{\text{WA}} = \frac{1}{K} \sum_{t=1}^K \theta_t
\end{equation}
Task Arithmetic (TA) scales and adds the task-specific updates back to the progenitor:
\begin{equation}
    \theta_{\text{TA}} = \theta_{\text{init}} + \lambda \sum_{t=1}^K \tau_t
\end{equation}
where $\lambda$ is a global scaling factor (typically $0.3$ to $0.5$ for $K=3$).

\subsection{State-of-the-Art Baselines: TIES-Merging and DARE}
To place our work in the context of recent developments, we include two highly influential modern parameter-space merging techniques:
\begin{enumerate}
    \item \textbf{TIES-Merging (Trimming, Sign Election, and Aggregation):} Yadav et al. \cite{yadav2023ties} address parameter interference by trimming each task update $\tau_t$ to keep only the top-$k\%$ highest-magnitude parameters (setting the rest to zero). A consensus sign vector $S = \text{sign}(\sum_t \text{sign}(\tau^{\text{trimmed}}_t))$ is then elected. Finally, the updates are aggregated by averaging only those parameters whose trimmed signs agree with the consensus sign.
    \item \textbf{DARE (Drop and Rescale):} Yu et al. \cite{yu2023language} apply a randomized dropout mask to the task updates. Elements in each update vector $\tau_t$ are randomly dropped (set to zero) with probability $p$, and the remaining elements are scaled up by $1/(1-p)$ to maintain the expected parameter update scale. The rescaled updates are then averaged.
\end{enumerate}
These techniques have shown strong results on large language models but have not been thoroughly tested under post-training quantization and data-efficient activation calibration in vision domains.

\subsection{Wasserstein-Calibrated Parameter Resonance (WCPR)}
WCPR \cite{theorist2026} treats task updates as continuous empirical distributions. It uses 1D Optimal Transport (OT) to align the merged task update $T_{\text{merged}} = \sum_t \tau_t$ channel-by-channel with the Wasserstein-2 barycenter of the expert updates. 

For each channel/row $c$ of a weight tensor $W$:
\begin{enumerate}
    \item Sort each expert's task update channel $\tau_{t, c}$ to obtain $\tau^{\text{sorted}}_{t, c}$.
    \item Compute the target barycenter: $T^{\text{sorted}}_{\text{target}, c} = \frac{1}{K} \sum_t \tau^{\text{sorted}}_{t, c}$.
    \item Replace the sorted elements of the merged task update $T_{\text{merged}, c}$ with $T^{\text{sorted}}_{\text{target}, c}$, preserving original ranks:
    \begin{equation}
        T^{\text{cal}, c}[i] = T^{\text{sorted}}_{\text{target}, c}[\text{rank}(T_{\text{merged}, c}[i])]
    \end{equation}
\end{enumerate}

\subsection{Proposed Quantization-Robust WCPR (QR-WCPR)}
While WCPR preserves exact higher-order statistics in FP32, its monotonic mapping can inflate the dynamic range of individual channels, mapping outlier elements to large values. Under post-training quantization, symmetric uniform quantization uses the maximum absolute value of a tensor to compute the step size $\Delta = \max(|W|) / q_{\text{max}}$. Outlier scale inflation thus adds severe rounding noise to all other weights in that channel.

To address this, we propose **Quantization-Robust WCPR (QR-WCPR)**. After performing 1D OT alignment to get $T^{\text{cal}, c}$, we calculate the maximum absolute update of each expert: $E^{\text{max}}_t = \| \tau_{t, c} \|_{\infty}$. We compute a robust target maximum absolute value using the median and Median Absolute Deviation (MAD) of the experts' maximum absolute values:
\begin{equation}
    M_c = \text{median}(E^{\text{max}}_1, \dots, E^{\text{max}}_K)
\end{equation}
\begin{equation}
    \text{MAD}_c = \text{median}(|E^{\text{max}}_t - M_c|)
\end{equation}
We define the upper bound on the channel's dynamic range as:
\begin{equation}
    max\_allowed = M_c + \gamma \cdot \text{MAD}_c
\end{equation}
If the maximum of the calibrated update exceeds $max\_allowed$, we scale the entire channel's update down:
\begin{equation}
    T^{\text{QR-WCPR}, c} = T^{\text{cal}, c} \cdot \min\left(1.0, \frac{max\_allowed}{\| T^{\text{cal}, c} \|_{\infty}}\right)
\end{equation}
This bounds the dynamic range of each weight channel, minimizing PTQ noise while preserving the rank and distribution shape of the optimal transport alignment.

\subsection{Data-Efficient BatchNorm Calibration (DE-BN)}
In networks with BatchNorm, the running statistics (mean and variance) of the experts become misaligned when weights are merged. Following \cite{methodologist2026}, we use a data-efficient calibration (DE-BN). We freeze all model weights and run a single forward pass with a small calibration subset of $N$ samples (e.g., $N=16$) in training mode. This estimates accurate cumulative moving averages for the BatchNorm layers, completely data-efficiently.

\section{Experimental Setup}
We specialized progenitor models on MNIST, Fashion-MNIST (FMNIST), and CIFAR-10. All datasets were resized to 3x32x32. We evaluated two architectures:
\begin{enumerate}
    \item **ResNet-18:** A standard residual network \cite{he2016deep} containing BatchNorm layers \cite{ioffe2015batch}.
    \item **MLP:** A deep Multi-Layer Perceptron with 4 hidden layers of size 512 and ReLU activations. This architecture contains no normalization layers.
\end{enumerate}

We evaluated merged models under nine distinct scenarios spanning full-precision (FP32), symmetric channel-wise post-training quantization to INT8 and INT4, and out-of-distribution environmental corruptions (Gaussian noise and Gaussian blur with severity 1.5).

\section{Results and Analysis}

Our empirical findings are summarized in Tables 1 to 4.

\begin{table*}[t]
\caption{Multi-task model merging accuracy (\%) in FP32 Clean across 3 random seeds (mean $\pm$ std).}
\label{tab:fp32_clean}
\begin{center}
\begin{tabular}{lcccc}
\toprule
\textbf{Architecture \& Method} & \textbf{MNIST} & \textbf{FMNIST} & \textbf{CIFAR-10} & \textbf{Average Acc} \\
\midrule
\textbf{ResNet-18 (Oracles)} & 98.91 & 92.39 & 78.90 & 90.07 \\
\midrule
ResNet-18 WA (No Calib) & 13.44 & 19.57 & 15.22 & 16.08 \\
ResNet-18 TA ($\lambda=0.3$, No Calib) & 11.54 & 17.54 & 15.26 & 14.78 \\
ResNet-18 TA ($\lambda=0.5$, No Calib) & 9.80 & 10.00 & 10.00 & 9.93 \\
ResNet-18 QR-IPR & 9.80 & 10.00 & 10.00 & 9.93 \\
ResNet-18 TIES ($\lambda=0.5$, No Calib) & 8.92 & 10.39 & 10.32 & 9.88 \\
ResNet-18 DARE ($\lambda=0.5$, No Calib) & 9.80 & 10.00 & 10.00 & 9.93 \\
ResNet-18 WCPR & 16.76 & 20.13 & 10.92 & 15.94 \\
ResNet-18 QR-WCPR & 15.60 & 20.84 & 11.38 & 15.94 \\
\midrule
ResNet-18 WA + DE-BN (16) & $47.20 \pm 0.86$ & $42.91 \pm 2.10$ & $38.37 \pm 1.82$ & $42.83 \pm 1.11$ \\
ResNet-18 TA ($\lambda=0.3$) + DE-BN (16) & $49.79 \pm 0.63$ & $45.02 \pm 1.57$ & $37.35 \pm 1.66$ & $44.05 \pm 0.96$ \\
ResNet-18 TA ($\lambda=0.5$) + DE-BN (16) & $32.14 \pm 0.16$ & $31.96 \pm 1.30$ & $\textbf{41.87} \pm 2.27$ & $35.32 \pm 1.06$ \\
ResNet-18 QR-IPR + DE-BN (16) & $33.37 \pm 0.46$ & $33.99 \pm 1.36$ & $41.71 \pm 2.13$ & $36.36 \pm 1.00$ \\
ResNet-18 TIES ($\lambda=0.5$) + DE-BN (16) & $\textbf{64.23} \pm 0.26$ & $\textbf{50.03} \pm 0.43$ & $30.10 \pm 1.45$ & $\textbf{48.12} \pm 0.41$ \\
ResNet-18 DARE ($\lambda=0.5$) + DE-BN (16) & $32.55 \pm 0.58$ & $33.05 \pm 1.25$ & $39.97 \pm 2.58$ & $35.19 \pm 0.97$ \\
ResNet-18 WCPR + DE-BN (16) & $42.61 \pm 0.66$ & $37.39 \pm 1.34$ & $41.04 \pm 2.03$ & $40.35 \pm 0.86$ \\
ResNet-18 QR-WCPR + DE-BN (16) & $43.07 \pm 0.64$ & $37.75 \pm 1.29$ & $40.67 \pm 1.91$ & $40.50 \pm 0.81$ \\
\midrule
\midrule
\textbf{MLP (Oracles)} & 97.47 & 87.23 & 54.53 & 79.74 \\
\midrule
MLP WA & 28.68 & 30.61 & 23.65 & 27.65 \\
MLP TA ($\lambda=0.3$) & 27.81 & 30.02 & 23.21 & 27.01 \\
MLP TA ($\lambda=0.5$) & 31.47 & 33.76 & 23.61 & 29.61 \\
MLP QR-IPR & 32.07 & 33.71 & 23.73 & 29.84 \\
MLP TIES ($\lambda=0.5$) & 37.96 & 24.61 & 18.19 & 26.92 \\
MLP DARE ($\lambda=0.5$) & 30.72 & 32.06 & 22.84 & 28.54 \\
MLP WCPR & 39.27 & $\textbf{42.25}$ & $\textbf{29.62}$ & $\textbf{37.05}$ \\
MLP QR-WCPR & $\textbf{39.74}$ & 41.70 & 29.37 & 36.94 \\
\bottomrule
\end{tabular}
\end{center}
\end{table*}

\begin{figure*}[t]
\centering
\includegraphics[width=0.95\textwidth]{plot_de_bn_sweep.png}
\caption{The data-efficiency curve of activation-space calibration (DE-BN) on ResNet-18 under clean FP32 (left) and extreme INT4 quantization (right). While weight-space calibration alone is completely ineffective (0 calibration samples), calibrating with as few as 1 to 16 samples dramatically heals representation collapse across all methods. Standard deviation is shaded across 3 random seeds.}
\label{fig:de_bn_sweep}
\end{figure*}

\subsection{Verification of the BatchNorm Gatekeeper Hypothesis}
As demonstrated in Table 1, under FP32 Clean, every single merging technique collapses completely (obtaining 9\% to 16\% accuracy, which is equivalent to random guessing) when evaluated on ResNet-18 without BatchNorm calibration. Weight-space calibration (WCPR, QR-IPR) does not solve this collapse.

However, when we activate DE-BN with only 16 samples, the collapse is resolved across all methods: Weight Averaging rises to $\textbf{42.83}\%$ ($\text{std}=1.11\%$) and Task Arithmetic ($\lambda=0.3$) rises to $\textbf{44.05}\%$ ($\text{std}=0.96\%$). WCPR + DE-BN obtains $\textbf{40.35}\%$ ($\text{std}=0.86\%$), and our proposed QR-WCPR + DE-BN obtains $\textbf{40.50}\%$ ($\text{std}=0.81\%$). This provides overwhelming empirical proof that BatchNorm statistics misalignment is the primary gatekeeper of representation collapse in normalization-based networks.

This data-efficiency behavior is mapped out in detail in Figure~\ref{fig:de_bn_sweep}. We sweep the number of calibration samples from 0 to 256 across 3 random seeds. Across all merging techniques, we observe a sharp, monotonic increase in performance as soon as even a single sample is used for calibration, with performance largely plateauing beyond 16 to 32 samples. Under both clean FP32 and INT4 quantization regimes, our findings show that extremely light-weight calibration (e.g., 16 samples, representing less than 0.03\% of the training data) is sufficient to completely restore activation statistics and heal representation collapse.

\subsection{WCPR on BatchNorm-less Architectures}
In our BatchNorm-less MLP architecture, where activation statistics are not constrained, weight-space calibration alone is highly effective. As shown in Table 1, standard WCPR achieves **37.05\%** average accuracy, outperforming Weight Averaging by **+9.4\%** (37.05\% vs 27.65\%) and Task Arithmetic ($\lambda=0.5$) by **+7.44\%**. This empirically validates the theoretical advantage of Wasserstein barycenter alignment in recovering joint representations.

\begin{table}[h]
\caption{Multi-task accuracy (\%) under INT8 Quantization across 3 random seeds (mean $\pm$ std).}
\label{tab:int8}
\begin{center}
\footnotesize
\setlength{\tabcolsep}{3pt}
\begin{tabular}{lcc}
\toprule
\textbf{Method} & \textbf{No BN Calib} & \textbf{DE-BN (16)} \\
\midrule
\textbf{ResNet-18 INT8} & & \\
Weight Averaging & 16.09 & $42.70 \pm 1.08$ \\
Task Arithmetic ($\lambda=0.3$) & 14.77 & $44.00 \pm 0.96$ \\
Task Arithmetic ($\lambda=0.5$) & 9.93 & $35.23 \pm 0.99$ \\
QR-IPR & 9.93 & $36.51 \pm 0.99$ \\
TIES ($\lambda=0.5$) & 9.91 & $\textbf{47.96} \pm 0.43$ \\
DARE ($\lambda=0.5$) & 9.93 & $35.32 \pm 0.96$ \\
WCPR & 15.91 & $40.18 \pm 0.91$ \\
QR-WCPR (Proposed) & 15.86 & $40.29 \pm 0.80$ \\
\midrule
\textbf{MLP INT8} & & \\
Weight Averaging & 27.62 & - \\
Task Arithmetic ($\lambda=0.3$) & 27.07 & - \\
Task Arithmetic ($\lambda=0.5$) & 29.59 & - \\
QR-IPR & 29.92 & - \\
TIES ($\lambda=0.5$) & 26.94 & - \\
DARE ($\lambda=0.5$) & 28.53 & - \\
WCPR & \textbf{37.11} & - \\
QR-WCPR (Proposed) & 37.01 & - \\
\bottomrule
\end{tabular}
\end{center}
\end{table}

\begin{figure}[t]
\centering
\includegraphics[width=0.48\textwidth]{plot_gamma_sweep.png}
\caption{Sensitivity analysis of the robust clamping threshold $\gamma$ under extreme INT4 quantization for ResNet-18 and MLP architectures. Bounding the dynamic range of Optimal Transport weight updates ($\gamma \approx 1.5$) successfully stabilizes quantization step size, outperforming the standard WCPR baseline (dashed lines).}
\label{fig:gamma_sweep}
\end{figure}

\subsection{PTQ Robustness and Outlier-Clamping (INT8 and INT4)}
Under 8-bit quantization (Table 2), both WCPR and QR-WCPR remain remarkably stable on both architectures. WCPR + DE-BN drops only slightly from $40.35\%$ to $40.18\%$ ($\text{std}=0.91\%$), while our proposed QR-WCPR + DE-BN drops from $40.50\%$ to $40.29\%$ ($\text{std}=0.80\%$).

We conduct a sensitivity analysis of the robust clamping threshold $\gamma$ under extreme INT4 quantization in Figure~\ref{fig:gamma_sweep}. When $\gamma$ is too small (e.g., 0.1), excessive clamping distorts the merged weight representations, leading to suboptimal accuracy. Conversely, when $\gamma$ is too large (e.g., 10.0), QR-WCPR behaves identically to standard WCPR, leaving the model vulnerable to outlier weight channels that inflate the dynamic range and amplify quantization noise. We find that a robust multiplier of $\gamma \in [1.0, 2.0]$ provides the optimal trade-off across both ResNet-18 and MLP architectures, successfully mitigating representation degradation and consistently outperforming the standard uncalibrated WCPR baselines.

\begin{table}[h]
\caption{Multi-task accuracy (\%) under INT4 Quantization across 3 random seeds (mean $\pm$ std).}
\label{tab:int4}
\begin{center}
\footnotesize
\setlength{\tabcolsep}{3pt}
\begin{tabular}{lcc}
\toprule
\textbf{Method} & \textbf{No BN Calib} & \textbf{DE-BN (16)} \\
\midrule
\textbf{ResNet-18 INT4} & & \\
Weight Averaging & 13.59 & $34.73 \pm 1.37$ \\
Task Arithmetic ($\lambda=0.3$) & 13.20 & $37.59 \pm 0.82$ \\
Task Arithmetic ($\lambda=0.5$) & 9.93 & $32.18 \pm 1.02$ \\
QR-IPR & 9.93 & $33.52 \pm 1.01$ \\
TIES ($\lambda=0.5$) & 10.69 & $\textbf{42.50} \pm 0.11$ \\
DARE ($\lambda=0.5$) & 9.93 & $32.53 \pm 1.06$ \\
WCPR & 16.55 & $33.39 \pm 1.10$ \\
QR-WCPR (Proposed) & 16.66 & $33.84 \pm 1.25$ \\
\midrule
\textbf{MLP INT4} & & \\
Weight Averaging & 27.71 & - \\
Task Arithmetic ($\lambda=0.3$) & 27.28 & - \\
Task Arithmetic ($\lambda=0.5$) & 29.16 & - \\
QR-IPR & 29.88 & - \\
TIES ($\lambda=0.5$) & 26.99 & - \\
DARE ($\lambda=0.5$) & 28.29 & - \\
WCPR & \textbf{36.73} & - \\
QR-WCPR (Proposed) & 36.42 & - \\
\bottomrule
\end{tabular}
\end{center}
\end{table}

However, under extreme 4-bit quantization (Table 3), the precision of parameters is highly constrained. For ResNet-18, standard WCPR + DE-BN drops to $\textbf{33.39}\%$ ($\text{std}=1.10\%$). Our proposed $\textbf{QR-WCPR + DE-BN}$ successfully mitigates this degradation, achieving $\textbf{33.84}\%$ ($\text{std}=1.25\%$) and outperforming standard WCPR. This confirms our hypothesis: by statistically bounding the channel-wise dynamic range of Optimal Transport weight updates, QR-WCPR stabilizes the quantization step size and improves robustness under extreme parameter precision limits.

\begin{table*}[t]
\caption{ResNet-18 accuracy (\%) under environmental corruptions (severity 1.5) across 3 random seeds (mean $\pm$ std). ``Uncal.'' denotes no activation calibration, while ``Cal.'' denotes DE-BN (16) calibration.}
\label{tab:corruptions}
\begin{center}
\footnotesize
\setlength{\tabcolsep}{4pt}
\begin{tabular}{lcccc}
\toprule
\textbf{Method} & \textbf{Noise (Uncal.)} & \textbf{Noise (Cal.)} & \textbf{Blur (Uncal.)} & \textbf{Blur (Cal.)} \\
\midrule
Weight Averaging & $16.85 \pm 0.03$ & $25.89 \pm 0.14$ & 15.75 & $39.10 \pm 1.57$ \\
Task Arithmetic ($\lambda=0.3$) & $16.65 \pm 0.04$ & $26.09 \pm 0.10$ & 14.17 & $40.50 \pm 1.53$ \\
Task Arithmetic ($\lambda=0.5$) & 9.93 & $25.78 \pm 0.52$ & 9.93 & $30.92 \pm 1.45$ \\
QR-IPR & 9.93 & $25.38 \pm 0.56$ & 9.93 & $32.30 \pm 1.46$ \\
TIES ($\lambda=0.5$) & $10.18 \pm 0.04$ & $\textbf{26.33} \pm 0.34$ & 9.88 & $\textbf{44.92} \pm 1.00$ \\
DARE ($\lambda=0.5$) & 9.93 & $25.47 \pm 0.60$ & 9.93 & $31.15 \pm 1.19$ \\
WCPR & $13.55 \pm 0.02$ & $25.44 \pm 0.68$ & \textbf{15.75} & $35.95 \pm 1.32$ \\
QR-WCPR (Proposed) & $\textbf{14.99} \pm 0.05$ & $25.44 \pm 0.70$ & 15.38 & $36.21 \pm 1.24$ \\
\bottomrule
\end{tabular}
\end{center}
\end{table*}

\subsection{Environmental Noise and Blur Robustness}
Under Gaussian Noise and Blur corruptions (Table 4), the models show a consistent and insightful trend across both architectures. Without DE-BN on ResNet-18, both Gaussian Noise and severe Gaussian Blur corruptions result in complete representation collapse (accuracies around 14\%--16\% for noise, and 9\%--15\% for blur, equivalent to random guessing). 

Activating DE-BN with only 16 samples dramatically rescues all merging techniques: raising Gaussian Noise performance to around 25\%--26\% (with Task Arithmetic $\lambda=0.3$ achieving $\textbf{26.09}\% \pm 0.10\%$), and raising Gaussian Blur performance to around 36\%--40\% (with Task Arithmetic achieving $\textbf{40.50}\% \pm 1.53\%$ and our proposed QR-WCPR achieving $\textbf{36.21}\% \pm 1.24\%$). 

Furthermore, even without activation calibration, our proposed QR-WCPR outperforms standard uncalibrated WCPR under Gaussian Noise ($14.99\% \pm 0.05\%$ vs $13.55\% \pm 0.02\%$), showing that robust dynamic range bounding preserves out-of-distribution (OOD) generalization.

For our BatchNorm-less MLP architecture, where activation stats are unconstrained, we observe excellent OOD generalization in full precision without any calibration. Under Gaussian Noise, standard WCPR achieves $36.57\% \pm 0.07\%$ and QR-WCPR achieves $36.44\% \pm 0.05\%$, vastly outperforming uncalibrated Weight Averaging ($26.65\% \pm 0.03\%$). Under Gaussian Blur, WCPR achieves $35.78\%$ and QR-WCPR achieves $35.72\%$, once again outperforming Weight Averaging ($26.01\%$). This confirms that Optimal Transport parameter-space mapping yields structural weight representations that are naturally robust to input-space environmental corruptions.

\subsection{Sensitivity Analyses: Merging Coefficient and PTQ Bitwidth}
In this section, we present additional extensive sweeps exploring the sensitivity of merged models to two crucial parameters: the merging scaling coefficient $\lambda$ and the physical quantization bitwidth.

\begin{figure*}[t]
\centering
\begin{subfigure}{0.49\textwidth}
    \centering
    \includegraphics[width=\textwidth]{plot_lambda_sweep.png}
    \caption{Merging Coefficient $\lambda$ Sensitivity}
    \label{fig:lambda_sweep}
\end{subfigure}
\hfill
\begin{subfigure}{0.49\textwidth}
    \centering
    \includegraphics[width=\textwidth]{plot_bitwidth_sweep.png}
    \caption{PTQ Bitwidth Robustness}
    \label{fig:bitwidth_sweep}
\end{subfigure}
\caption{Comprehensive sweeps of merging scaling coefficient and physical quantization bitwidth. (a) Scaling coefficient $\lambda$ sweep under FP32 Clean (with DE-BN 16 for ResNet-18). (b) Quantization precision sweep from 2-bit to 8-bit. Standard WCPR and our proposed QR-WCPR exhibit complete scaling invariance, eliminating a major hyperparameter tuning overhead.}
\label{fig:additional_sweeps}
\end{figure*}

\paragraph{Scaling Invariance of 1D Optimal Transport Merging}
As illustrated in Figure~\ref{fig:lambda_sweep}, we sweep the merging coefficient $\lambda$ from 0.1 to 1.0. For standard Task Arithmetic, the choice of $\lambda$ is extremely sensitive and fragile. On ResNet-18, TA peaks at $\lambda=0.1$ with 47.81\% accuracy and rapidly collapses to 22.51\% at $\lambda=1.0$. On MLP, TA peaks at $\lambda=0.7$ with 30.18\% and collapses to 22.09\% at $\lambda=0.1$. This requires costly, task-specific hyperparameter search.

In stark contrast, both WCPR and our proposed QR-WCPR exhibit complete invariance to the scaling coefficient $\lambda$. Under all tested $\lambda \in [0.1, 1.0]$, WCPR and QR-WCPR maintain exactly constant accuracy (40.71\% and 40.94\% on ResNet-18, and 37.05\% and 36.94\% on MLP). 

This is a profound mathematical property of 1D Optimal Transport merging. The OT map in Equation (3) operates on the channel-wise updates using the sorting/argsort operator. Because any positive scalar multiplication $x \leftarrow \lambda \cdot x$ preserves the strict rank order of the tensor elements (i.e., $\text{argsort}(\lambda \cdot x) = \text{argsort}(x)$ for any $\lambda > 0$), the rearrangement indices are identical. Thus, the mapping to the target Wasserstein barycenter is completely invariant to the magnitude of the update. This eliminates the need for any lambda search, greatly simplifying deploying model merging.

\paragraph{Bitwidth Breakdown Curve}
In Figure~\ref{fig:bitwidth_sweep}, we sweep the physical post-training quantization bitwidth from 2-bit to 8-bit. Under extreme 2-bit quantization, all methods collapse to near random guessing ($\sim$10\%-12\%) due to the severe loss of precision. However, starting from 3-bit, we observe distinct performance profiles.

On ResNet-18, our proposed QR-WCPR consistently outperforms standard WCPR across low bitwidths: achieving 28.79\% vs 28.17\% in 3-bit, 34.59\% vs 33.83\% in 4-bit, and 40.03\% vs 39.37\% in 5-bit, validating that robust clamping of outlier channels significantly mitigates PTQ rounding noise. On MLP, the gap is also clearly visible in 3-bit (33.84\% for QR-WCPR vs 32.98\% for WCPR). Both WCPR variants consistently and significantly outperform uncalibrated Weight Averaging and Task Arithmetic at 4-bit, 5-bit, 6-bit, and 8-bit, proving that Wasserstein barycenter alignment is highly resilient to physical quantization down to extreme limits.

\section{Discussion and Related Work}
Our findings bridge the gap between several key directions in multi-task model merging, post-training quantization, and deep learning optimization.

\paragraph{Multi-Task Model Merging}
Multi-task merging has emerged as a promising training-free paradigm to consolidate specialized networks into a single, cohesive model \cite{ilharco2023editing, wortsman2022robust, wortsman2022model, choshen2022fusing, ortiz2023task}. Uncalibrated merging, however, causes destructive interference and representation collapse. Early solutions integrated permutations or alignments modulo symmetries \cite{ainsworth2023git, stoica2023zipit, NEURIPS2020_Singh}, computed Fisher information weights \cite{matena2022merging}, or pruned updates to mitigate task interference \cite{yadav2023ties, yu2023language}. While parameter-level transport like WCPR \cite{theorist2026} maps weights to Wasserstein-2 barycenters, our work reveals a fundamental limitation: purely parameter-space alignments fail in normalization-based networks (such as ResNets \cite{he2016deep}) because of activation statistics misalignment. Conversely, while activation calibration like DE-BN \cite{methodologist2026} resolves this misalignment, it does not correct structural weight-space distortion. This highlights the necessity of co-designing weight alignment with activation-level calibration.

\paragraph{Post-Training Quantization (PTQ)}
PTQ is critical for deploying deep neural networks under strict hardware and edge constraints \cite{han2015deep, jacob2018quantization}. Landmark approaches like GPTQ \cite{frantar2023gptq}, QLoRA \cite{dettmers2023qlora}, AdaRound \cite{nagel2020up}, and general surveys \cite{gholami2021survey, krishnamoorthi2018quantizing} optimize rounding noise and step sizes. In the context of model merging, however, alignment techniques can inflate the dynamic range of individual channels. Our proposed QR-WCPR resolves this by introducing robust statistical outlier clamping, preserving Wasserstein alignment while stabilizing the physical quantization scale on low-bit hardware (down to 3-bit and 4-bit).

\paragraph{Symmetries, Loss Surfaces, and Optimization Theory}
The study of loss surfaces, mode connectivity, and implicit optimization bias provides the mathematical foundation of weight aggregation \cite{choromanska2015loss, keskar2016on, hochreiter1997flat}. Research in the lottery ticket hypothesis and linear mode connectivity \cite{frankle2019lottery, frankle2020linear, neyshabur2020what} demonstrates that models starting from a shared initialization remain in the same loss basin. Our study of scaling invariance in WCPR and QR-WCPR is deeply rooted in optimization and statistical learning theory \cite{vapnik1998statistical, bartlett2002rademacher, shalev2014understanding, soudry2018implicit, gunasekar2018characterizing, neyshabur2015search, jacot2018neural, arora2019fine}, showing how sorting operations preserve rank representation under positive scalar transformations. These insights are highly relevant as model merging scales to modern large-scale architectures, including Transformers \cite{vaswani2017attention}, BERT \cite{devlin2018bert}, GPT \cite{radford2019language, brown2020language, achiam2023gpt}, and LLaMA \cite{touvron2023llama}, where training from scratch is computationally prohibitive \cite{zhao2023survey}.

\subsection{Limitations and Future Directions}
While our multi-scenario empirical evaluation is extensive and offers key insights, we acknowledge several limitations that pave the way for future research:
\begin{itemize}
    \item \textbf{Model Architecture and Scale:} We focus on ResNet-18 and MLP networks on MNIST, Fashion-MNIST, and CIFAR-10. Although these settings are representative of core architectural components (with and without normalization) and allow us to run extensive, multi-seed sweeps, future research should investigate the scaling properties of WCPR, QR-WCPR, and DE-BN on larger-scale Transformer-based architectures, such as Vision Transformers (ViTs) and Large Language Models (LLMs).
    \item \textbf{Quantization Schemes:} Our post-training quantization is restricted to uniform, symmetric channel-wise quantization. Testing asymmetric quantization or non-uniform data types (e.g., NormalFloat4 or integer-only arithmetic) could reveal additional nuances in Optimal Transport weight-space stability.
    \item \textbf{Number of Experts:} We evaluate multi-task model merging under three expert networks ($K=3$). Scaling the number of experts to dozens or hundreds of tasks (as in continual learning or federated multi-task scenarios) might introduce larger distribution drifts that require dynamic clamping schedules.
\end{itemize}

\section{Conclusion}
In this work, we presented a massive, multi-scenario empirical evaluation of data-free and data-efficient model merging. We demonstrated that BatchNorm running statistics misalignment is the primary cause of representation collapse in ResNet-18, which is completely healed by data-efficient BatchNorm calibration (DE-BN) using only 16 samples. Conversely, in MLP architectures lacking BatchNorm, 1D Optimal Transport (WCPR) provides a major, robust data-free gain of +9.4\% accuracy. Finally, under extreme 4-bit quantization, our proposed QR-WCPR leverages robust statistical clamping to stabilize weight dynamic ranges, mitigating quantization noise and outperforming standard WCPR. Our work highlights the critical need to co-design parameter alignment and activation calibration for robust edge deployment.

\bibliography{example_paper}
\bibliographystyle{icml2026}

\end{document}
